% !PRINTED 423
\weberpage{423}

%<<CONT
\weq{2}{Y^5 + 5aY^2 + 5bY + c = 0,}
in which $a, b, c$ denote homogeneous functions of degrees 3, 4, 5 in $\alpha, \beta$, whose coefficients are icosahedral invariants.

We can easily determine the function $c$ from the formula
\[ c = -\Pi(\tilde\omega)\,\Pi(\alpha + \beta t). \]

For, by (15) of the preceding section, we have for an indeterminate $\lambda$
\[ \lambda^5 - 10f\lambda^3 + 45f^2\lambda - T = \Pi(\lambda - t), \]
hence, if one sets $\lambda = -\alpha : \beta$,
\[ \Pi(\alpha + \beta t) = \alpha^5 - 10f\alpha^3\beta^2 + 45f^2\alpha\beta^4 + T\beta^5, \]
and further, by comparison of a single term, there follows at once [\S.\,111, (6), (8)]
\[ \Pi(\tilde\omega) = -H^2, \]
hence
\weq{3}{c = H^2(\alpha^5 - 10f\alpha^3\beta^2 + 45f^2\alpha\beta^4 + T\beta^5).}

The coefficients $a, b$ are perhaps most simply computed with the help of the \gesperrt{Newton} formulae (Bd.\,I, \S.\,42):
\[
\begin{aligned}
-15a &= S(\alpha\tilde\omega + \beta t\tilde\omega)^3\\
-20b &= S(\alpha\tilde\omega + \beta t\tilde\omega)^4,
\end{aligned}
\]
by applying the formulae which result, after a short calculation, from the comparison of one term each of the right and left sides:
\[
\begin{array}{ll}
S(\tilde\omega^3) = -3.5.8f^2, & S(t\tilde\omega^3) = -5T,\\
S(t^2\tilde\omega^3) = -5.72f^3, & S(t^3\tilde\omega^3) = -3.5fT,\\
S(\tilde\omega^4) = 4.5fH, & S(t^2\tilde\omega^4) = -5.12f^2H,\\
S(t^3\tilde\omega^4) = -5HT, & S(t^4\tilde\omega^4) = -4.5.27f^3H.
\end{array}
\]

Thus one finds
\weq{4}{a = 8f^2\alpha^3 + T\alpha^2\beta + 72f^3\alpha\beta^2 + fT\beta^3}
\weq{5}{b = -fH\alpha^4 + 18f^2H\alpha^2\beta^2 + HT\alpha\beta^3 + 27f^3H\beta^4.}

In order to obtain from this the resolvent of the icosahedral equation, we must introduce in place of $\tilde\omega$ and $t$ the functions $u, v$ [\S.\,111, (12)].

If we set, denoting by $\lambda, \mu$ any two indeterminate quantities,
\weq{6}{\alpha = \frac{\lambda f}{H},\quad \beta = \frac{\mu f^3}{HT},}
%CONT>>
